% AdvancedEMQuality20261023
\section{量子力学初步}

\subsection{Schrödinger 方程与定态问题}

把经典 Hamilton 量中的
$\vect p$ 换成算符 $-\ii\hbar\nabla$，得到
\[
 \ii\hbar\frac{\partial\psi}{\partial t}
 =\widehat H(\vect r,-\ii\hbar\nabla)\psi.
\]
若势能不显含时间，可分离为
$\psi(\vect r,t)=\Psi(\vect r)\e^{-\ii Et/\hbar}$，空间波函数满足
\[
 \boxed{\widehat H\Psi=E\Psi}.
\]
自由粒子
$\widehat H=-\hbar^2\nabla^2/(2m)$
的平面波满足 $E=\hbar\omega$、$\vect p=\hbar\vect k$、
$E=p^2/(2m)$。

\begin{advexamplebox}{球形禁闭、动能下界与氢原子稳定性}\label{advA6:ex:hydrogen-confinement}
把电子波函数完全限制在半径为 $R$ 的球内并令边界上 $\psi=0$。

\begin{enumerate}
\item 证明动能存在下界，取下界的波函数满足
$\widehat T\psi=T_{\min}\psi$；
\item 求 $T_{\min}$ 与基态波函数；
\item 说明 Coulomb 势下电子不能无限塌缩到原点；
\item 用该球内基态作变分试探函数，估算氢原子基态能量。
\end{enumerate}

\begin{center}\begin{minipage}{.96\linewidth}\centering
\begin{tikzpicture}[>=Latex,font=\small,scale=.85,every node/.style={fill=none}]

\draw[thick] (0,0) circle (1.3);\fill (0,0) circle (1.5pt) node[below]{$O$};\draw[->] (0,0)--(1.3,0) node[midway,below]{$R$};\node[above] at (0,1.6){$\psi(R)=0$};\node[right] at (1.7,.5){球内允许，球外 $\psi=0$};

\end{tikzpicture}
\figcap{题图：球形禁闭区域与Dirichlet边界；Coulomb问题另以此态作变分试探。}\label{advA6:quality:A6-quantum-mechanics:q1}
\end{minipage}\end{center}
\end{advexamplebox}
\global\AdvHideHeadingtrue
\begin{advsolution}
\begin{center}\begin{minipage}{.96\linewidth}\centering
% EMAdvFigureHeading
\raggedright{\bfseries\color{DeepBlue}解答}\par\smallskip\centering

\begin{tikzpicture}[>=Latex,font=\small,scale=.85,every node/.style={fill=none}]

\draw[thick] (0,0) circle (1.3);\fill (0,0) circle (1.5pt) node[below]{$O$};\draw[->] (0,0)--(1.3,0) node[midway,below]{$R$};\node[above] at (0,1.6){$\psi(R)=0$};\node[right] at (1.7,.5){球内允许，球外 $\psi=0$};

\end{tikzpicture}
\figcap{解图：球形禁闭区域与Dirichlet边界；Coulomb问题另以此态作变分试探。}\label{advA6:quality:A6-quantum-mechanics:a1}
\end{minipage}\end{center}

归一化约束下令
\[
 \mathcal F[\psi]=
 \int_{r<R}\psi^*
 \left(-\frac{\hbar^2}{2m}\nabla^2\right)\psi\,\dd V
 -\lambda\left(\int|\psi|^2\dd V-1\right).
\]
对 $\psi^*$ 变分得
$\widehat T\psi=\lambda\psi$。因此最小值由 Dirichlet Laplacian 的最低本征值给出。

球对称基态令 $u=r\psi$，满足
$u''+k^2u=0$、$u(0)=u(R)=0$。故
\[
 k=\frac\pi R,\qquad
 \boxed{T_{\min}=\frac{\pi^2\hbar^2}{2mR^2}},
\]
\[
 \boxed{\psi_R(r)=
 \frac1{\sqrt{2\pi R}}\frac{\sin(\pi r/R)}r
 \quad(0<r<R)}.
\]
尺度压缩 $R\to0$ 时动能按 $R^{-2}$ 发散，而 Coulomb 势能只按
$-R^{-1}$ 发散，所以总能量必在有限尺度取最小值。

上述是尺度解释，不能只凭一族试探函数就断言所有态都有下界。对任意归一化、有限动能且使分部积分面项消失的波函数，令 $a_0=\hbar^2/(mk_ee^2)$，利用 $\nabla\cdot\hat{\boldsymbol r}=2/r$ 可直接配方：
\[
 \langle H\rangle=
 \frac{\hbar^2}{2m}\int\left|\nabla\psi+
 \frac{\hat{\boldsymbol r}}{a_0}\psi\right|^2\mathrm dV
 -\frac{\hbar^2}{2ma_0^2}
 \ge-\frac{mk_e^2e^4}{2\hbar^2}.
\]
这里 $k_e=1/(4\pi\varepsilon_0)$；平方项展开后的交叉项恰是 $-k_ee^2\langle1/r\rangle$。这个对所有允许态成立的界排除了无限塌缩。等号态的径向衰减为 $e^{-r/a_0}$，也说明下面的硬球试探态只能给出较粗的变分上界。


记 $k_e=1/(4\pi\varepsilon_0)$。对上述试探态
\[
 \left\langle\frac1r\right\rangle
 =\frac{C}{R},\qquad
 C=2\int_0^\pi\frac{\sin^2u}{u}\,\dd u
 =2.43765\ldots
\]
故
\[
 E(R)=\frac{\pi^2\hbar^2}{2mR^2}
 -\frac{Ck_ee^2}{R}.
\]
极小点与能量为
\[
 R_*=\frac{\pi^2}{C}a_0=4.049\,a_0,
\]
\[
 \boxed{E_{\rm var}=
 -\frac{C^2}{\pi^2}\frac{mk_e^2e^4}{2\hbar^2}
 =-8.19\,\mathrm{eV}}.
\]
它高于精确基态 $-13.6\,\mathrm{eV}$，符合变分上界原理；估算较粗是因硬球边界并非真实指数尾部。

图~\ref{advA6:quality:A6-quantum-mechanics:a1} 中，半径 $R$ 同时决定边界和径向驻波波数；球内基态用 $u=r\psi$ 消去球坐标的一阶导数，变分能量还须连同势能一起计算。
\end{advsolution}

\subsection{三个一维模型}

\begin{advexamplebox}{无限深方势阱}\label{advA6:ex:infinite-well}
粒子被限制在 $0<x<a$，阱外势能无穷大。证明能级和归一化波函数为
\[
 E_n=\frac{n^2\pi^2\hbar^2}{2ma^2},\qquad
 \Psi_n(x)=\sqrt{\frac2a}\sin\frac{n\pi x}{a}.
\]

\begin{center}\begin{minipage}{.96\linewidth}\centering
\begin{tikzpicture}[>=Latex,font=\small,scale=.85,every node/.style={fill=none}]

\draw[->] (-.7,0)--(4.2,0) node[right]{$x$};\draw[very thick] (0,2)--(0,0)--(3,0)--(3,2);\node[below] at (0,0){0};\node[below] at (3,0){$a$};\node at (1.5,1){$V=0$};\node[left] at (0,2){$\infty$};\node[right] at (3,2){$\infty$};

\end{tikzpicture}
\figcap{题图：无限深阱的两个硬壁位于 $0,a$，波函数在两壁为零。}\label{advA6:quality:A6-quantum-mechanics:q2}
\end{minipage}\end{center}
\end{advexamplebox}
\global\AdvHideHeadingtrue
\begin{advsolution}
\begin{center}\begin{minipage}{.96\linewidth}\centering
% EMAdvFigureHeading
\raggedright{\bfseries\color{DeepBlue}解答}\par\smallskip\centering

\begin{tikzpicture}[>=Latex,font=\small,scale=.85,every node/.style={fill=none}]

\foreach \shift/\n/\freq in {0/1/60,5/2/120}{\begin{scope}[xshift=\shift cm]
\draw[->] (-.3,0)--(3.5,0) node[right]{$x$};\draw[->] (0,-.9)--(0,1.1) node[above]{$\Psi_{\n}$};\draw[dashed] (3,-.9)--(3,1.1);\node[below] at (3,-.9){$a$};
\draw[DeepBlue,thick,domain=0:3,samples=80] plot(\x,{.7*sin(\freq*\x)});\end{scope}}

\end{tikzpicture}
\figcap{解图：$n=1,2$ 的驻波在两端均为零，第二态还在阱中点有节点；两图纵轴共同按波函数幅值取示意比例。}\label{advA6:quality:A6-quantum-mechanics:a2}
\end{minipage}\end{center}

阱内定态方程为
\[
 -\frac{\hbar^2}{2m}\Psi''=E\Psi,
\quad
 \Psi=A\sin kx+B\cos kx.
\]
边界 $\Psi(0)=0$ 给 $B=0$，$\Psi(a)=0$ 给
$ka=n\pi$，$n=1,2,\ldots$。由 $E=\hbar^2k^2/(2m)$ 得所求能级。
归一化
\[
 1=|A|^2\int_0^a\sin^2\frac{n\pi x}{a}\,\dd x
 =|A|^2\frac a2
\]
给 $A=\sqrt{2/a}$。$n=0$ 只给零波函数，不是物理态；
$E_n\propto a^{-2}$ 与禁闭动能尺度一致。

图~\ref{advA6:quality:A6-quantum-mechanics:a2} 中，两个端点同时为节点，所以只有整数个半波可放入宽度 $a$；$n=0$ 的曲线恒为零，不能归一化，因而不构成态。
\end{advsolution}

\begin{advexamplebox}{有限方势垒的隧穿系数}\label{advA6:ex:square-barrier}
势能
\[
 V(x)=\begin{cases}
 V_0,&0<x<a,\\0,&\text{其他},
 \end{cases}
\qquad 0<E<V_0.
\]
证明透射系数为
\[
 T=\left[
 1+\frac{V_0^2\sinh^2(\kappa a)}
 {4E(V_0-E)}\right]^{-1},
\qquad
 \kappa=\frac{\sqrt{2m(V_0-E)}}{\hbar}.
\]

\begin{center}\begin{minipage}{.96\linewidth}\centering
\begin{tikzpicture}[>=Latex,font=\small,scale=.85,every node/.style={fill=none}]

\draw[->] (-1.5,0)--(4.5,0) node[right]{$x$};\draw[thick] (-1.4,0)--(0,0)--(0,1.5)--(2.5,1.5)--(2.5,0)--(4.2,0);
\draw[dashed] (-1.3,.8)--(4,.8) node[right]{$E$};\node[above] at (1.2,1.5){$V_0$};\node[below] at (0,0){0};\node[below] at (2.5,0){$a$};\node[below] at (-.8,-.5){I};\node[below] at (1.2,-.5){II};\node[below] at (3.5,-.5){III};

\end{tikzpicture}
\figcap{题图：三区方势垒，能量位于阱外势能与垒顶之间。}\label{advA6:quality:A6-quantum-mechanics:q3}
\end{minipage}\end{center}
\end{advexamplebox}
\global\AdvHideHeadingtrue
\begin{advsolution}
\begin{center}\begin{minipage}{.96\linewidth}\centering
% EMAdvFigureHeading
\raggedright{\bfseries\color{DeepBlue}解答}\par\smallskip\centering

\begin{tikzpicture}[>=Latex,font=\small,scale=.85,every node/.style={fill=none}]

\draw[->] (-1.5,0)--(4.5,0) node[right]{$x$};\draw[thick] (-1.4,0)--(0,0)--(0,1.5)--(2.5,1.5)--(2.5,0)--(4.2,0);
\draw[dashed] (-1.3,.8)--(4,.8) node[right]{$E$};\node[above] at (1.2,1.5){$V_0$};\node[below] at (0,0){0};\node[below] at (2.5,0){$a$};\node[below] at (-.8,-.5){I};\node[below] at (1.2,-.5){II};\node[below] at (3.5,-.5){III};

\end{tikzpicture}
\figcap{解图：三区方势垒，能量位于阱外势能与垒顶之间。}\label{advA6:quality:A6-quantum-mechanics:a3}
\end{minipage}\end{center}

三区波函数取
\[
 \Psi_I=\e^{\ii kx}+r\e^{-\ii kx},\quad
 \Psi_{II}=A\e^{\kappa x}+B\e^{-\kappa x},\quad
 \Psi_{III}=t\e^{\ii kx},
\]
$k=\sqrt{2mE}/\hbar$。在 $x=0,a$ 连续匹配 $\Psi,\Psi'$。
为看清消元所依据的四个匹配条件，势垒内部的传播可写成
\[
 \begin{pmatrix}\Psi(a)\\\Psi'(a)\end{pmatrix}
 =\begin{pmatrix}
 \cosh\kappa a&\sinh(\kappa a)/\kappa\\
 \kappa\sinh\kappa a&\cosh\kappa a
 \end{pmatrix}
 \begin{pmatrix}\Psi(0)\\\Psi'(0)\end{pmatrix}.
\]
代入 $\Psi(0)=1+r$、$\Psi'(0)=\mathrm ik(1-r)$，以及 $\Psi(a)=te^{\mathrm ika}$、$\Psi'(a)=\mathrm ikte^{\mathrm ika}$，两式相减即可消去 $r$。整理后，
\[
 t=\frac{\e^{-\ii ka}}
 {\cosh\kappa a+
 \frac{\ii}{2}\left(\frac\kappa k-\frac k\kappa\right)
 \sinh\kappa a}.
\]
左右介质相同，流密度比就是 $T=|t|^2$。利用
\[
 \frac14\left(\frac\kappa k+\frac k\kappa\right)^2
 =\frac{V_0^2}{4E(V_0-E)}
\]
得到题中结果。高厚势垒 $\kappa a\gg1$ 时
\[
 T\simeq
 \frac{16E(V_0-E)}{V_0^2}\e^{-2\kappa a},
\]
显示指数隧穿；$a\to0$ 时 $T\to1$。分母中的 $4$ 不可遗漏。

图~\ref{advA6:quality:A6-quantum-mechanics:a3} 中，I、III区具有同一实波数 $k$，II区具有实衰减系数 $\kappa$；因此透射流比可直接取 $|t|^2$，而势垒内须同时保留两个指数项以满足两个界面的匹配。
\end{advsolution}

\begin{advexamplebox}{一维量子谐振子}\label{advA6:ex:harmonic-oscillator}
证明
\[
 E_n=\left(n+\frac12\right)\hbar\omega,
\]
\[
 \Psi_n(x)=
 \frac{1}{\sqrt{2^nn!}}
 \left(\frac{m\omega}{\pi\hbar}\right)^{1/4}
 H_n(\xi)\e^{-\xi^2/2},
\qquad
 \xi=\sqrt{\frac{m\omega}{\hbar}}x.
\]

\begin{center}\begin{minipage}{.96\linewidth}\centering
\begin{tikzpicture}[>=Latex,font=\small,scale=.85,every node/.style={fill=none}]

\draw[->] (-2,0)--(2.3,0) node[right]{$x$};\draw[->] (0,-.2)--(0,2.2) node[above]{$V$};\draw[DeepBlue,domain=-1.7:1.7,samples=70] plot(\x,{.65*\x*\x});\node[right] at (1.2,.8){$V=m\omega^2x^2/2$};

\end{tikzpicture}
\figcap{题图：一维谐势及坐标原点，束缚态需在左右无穷远可归一化。}\label{advA6:quality:A6-quantum-mechanics:q4}
\end{minipage}\end{center}
\end{advexamplebox}
\global\AdvHideHeadingtrue
\begin{advsolution}
\begin{center}\begin{minipage}{.96\linewidth}\centering
% EMAdvFigureHeading
\raggedright{\bfseries\color{DeepBlue}解答}\par\smallskip\centering

\begin{tikzpicture}[>=Latex,font=\small,scale=.85,every node/.style={fill=none}]

\draw[->] (-1.8,0)--(2,0) node[right]{$x$};\draw[DeepBlue,domain=-1.5:1.5] plot(\x,{.7*\x*\x});\foreach \y/\n in {.35/0,.85/1,1.35/2}{\draw[thick] (-1.5,\y)--(1.5,\y);\node[right] at (1.5,\y){$E_{\n}$};}\draw[->] (-1.7,.4)--(-1.7,.8) node[left]{$a^\dagger$};

\end{tikzpicture}
\figcap{解图：谐势内的等间隔能级，最低能级高于势能最低点。}\label{advA6:quality:A6-quantum-mechanics:a4}
\end{minipage}\end{center}

定义
\[
 \widehat a=
 \sqrt{\frac{m\omega}{2\hbar}}\widehat x
 +\frac{\ii}{\sqrt{2m\hbar\omega}}\widehat p,\qquad
 [\widehat a,\widehat a^\dagger]=1.
\]
Hamilton 量分解为
\[
 \widehat H=\hbar\omega
 \left(\widehat a^\dagger\widehat a+\frac12\right).
\]
数算符 $\widehat N=\widehat a^\dagger\widehat a$ 半正定。
若其最低本征值不是 $0$，连续施用 $\widehat a$ 会得到更低本征值，矛盾；
故 $\widehat a|0\rangle=0$，$E_0=\hbar\omega/2$。
再由
\[
 |n\rangle=\frac{(\widehat a^\dagger)^n}{\sqrt{n!}}|0\rangle
\]
得到 $E_n=(n+1/2)\hbar\omega$。

位置表象中 $\widehat a\Psi_0=0$ 给
$\Psi_0\propto\e^{-\xi^2/2}$；反复作用
$\widehat a^\dagger$ 并用 Rodrigues 公式
\[
 H_n(\xi)=(-1)^n\e^{\xi^2}
 \frac{\dd^n}{\dd\xi^n}\e^{-\xi^2}
\]
得到题中归一化波函数。$n$ 偶、奇时波函数宇称分别为偶、奇，
与对称势能相符。

图~\ref{advA6:quality:A6-quantum-mechanics:a4} 中，抛物势的两端都上升，允许态要同时满足两端衰减；引入升降算符后，基态被降算符湮灭，逐级升高能量而不改变相邻间隔 $\hbar\omega$。
\end{advsolution}

\subsection{Hilbert 空间与 Dirac 符号}

\begin{advdefinition}{Hilbert 空间}\label{advA6:def:hilbert}
Hilbert 空间 $\mathcal H$ 是带正定 Hermitian 内积的完备复线性空间。
内积对第二个变量线性、对第一个变量反线性：
\[
 \langle a|\lambda b+\mu c\rangle
 =\lambda\langle a|b\rangle+\mu\langle a|c\rangle,
\quad
 \langle b|a\rangle=\langle a|b\rangle^*.
\]
范数与距离为
\[
 \|a\|=\sqrt{\langle a|a\rangle},\qquad
 d(a,b)=\|a-b\|.
\]
完备性指每个 Cauchy 序列都收敛到空间内的元素。
\end{advdefinition}

\begin{advexamplebox}{Hilbert 空间中的 Cauchy--Schwarz 不等式}\label{advA6:ex:cauchy-schwarz}
证明
\[
 |\langle a|b\rangle|^2
 \le\langle a|a\rangle\langle b|b\rangle.
\]

\begin{center}\begin{minipage}{.96\linewidth}\centering
\begin{tikzpicture}[>=Latex,font=\small,scale=.85,every node/.style={fill=none}]

\draw[DeepBlue,thick,->] (0,0)--(2.7,0) node[right]{$|a\rangle$};\draw[DeepBlue,thick,->] (0,0)--(1.5,1.6) node[above]{$|b\rangle$};\fill (0,0) circle (1pt) node[below]{$0$};

\end{tikzpicture}
\figcap{题图：两个Hilbert空间向量；图示为它们张成子空间的示意，复相位仍由内积规定。}\label{advA6:quality:A6-quantum-mechanics:q5}
\end{minipage}\end{center}
\end{advexamplebox}
\global\AdvHideHeadingtrue
\begin{advsolution}
\begin{center}\begin{minipage}{.96\linewidth}\centering
% EMAdvFigureHeading
\raggedright{\bfseries\color{DeepBlue}解答}\par\smallskip\centering

\begin{tikzpicture}[>=Latex,font=\small,scale=.85,every node/.style={fill=none}]

\draw[DeepBlue,thick,->] (0,0)--(2.7,0) node[right]{$|a\rangle$};\draw[DeepBlue,thick,->] (0,0)--(1.5,1.6) node[above]{$|b\rangle$};\fill (0,0) circle (1pt) node[below]{$0$};
\draw[dashed] (1.5,0)--(1.5,1.6);\fill (1.5,0) circle (1.5pt) node[below]{$\lambda|a\rangle$};\node[right] at (1.5,.8){$|b\rangle-\lambda|a\rangle$};
\end{tikzpicture}
\figcap{解图：经适当相位选择后，投影与正交余量的实平面截面。}\label{advA6:quality:A6-quantum-mechanics:a5}
\end{minipage}\end{center}

按解图把 $b$ 分解为沿 $a$ 的投影和正交余量。若 $a=0$ 显然成立；
若 $a\ne0$，对任意复数 $\lambda$ 都有
\[
 0\le\|b-\lambda a\|^2.
\]
取 $\lambda=\langle a|b\rangle/\langle a|a\rangle$，余量与 $a$ 正交，且
\[
 0\le\langle b|b\rangle
 -\frac{|\langle a|b\rangle|^2}{\langle a|a\rangle}.
\]
乘以正数 $\langle a|a\rangle$ 即得不等式。等号要求余量为零，
即 $b=\lambda a$；连同零向量情形，等号当且仅当两个向量线性相关。

图~\ref{advA6:quality:A6-quantum-mechanics:a5} 中，把 $|b\rangle$ 分成沿 $|a\rangle$ 的投影与正交余量，余量范数非负给出不等式；等号要求余量为零，即两个向量线性相关。
\end{advsolution}

\begin{advexamplebox}{Hilbert 空间中的三角不等式}\label{advA6:ex:triangle}
证明
\[
 d(a,b)\le d(a,c)+d(c,b).
\]

\begin{center}\begin{minipage}{.96\linewidth}\centering
\begin{tikzpicture}[>=Latex,font=\small,scale=.85,every node/.style={fill=none}]

\coordinate (a) at (0,0);\coordinate (b) at (3,.3);\coordinate (c) at (1.3,1.7);\draw[thick] (a)--(b)--(c)--cycle;\node[left] at (a){$a$};\node[right] at (b){$b$};\node[above] at (c){$c$};\node[below] at (1.5,.1){$d(a,b)$};\node[left] at (.6,.9){$d(a,c)$};\node[right] at (2.2,1){$d(c,b)$};

\end{tikzpicture}
\figcap{题图：三个态向量在其仿射实平面中的距离三角形，距离由Hilbert范数定义。}\label{advA6:quality:A6-quantum-mechanics:q6}
\end{minipage}\end{center}
\end{advexamplebox}
\global\AdvHideHeadingtrue
\begin{advsolution}
\begin{center}\begin{minipage}{.96\linewidth}\centering
% EMAdvFigureHeading
\raggedright{\bfseries\color{DeepBlue}解答}\par\smallskip\centering

\begin{tikzpicture}[>=Latex,font=\small,scale=.85,every node/.style={fill=none}]

\coordinate (a) at (0,0);\coordinate (b) at (3,.3);\coordinate (c) at (1.3,1.7);\draw[thick] (a)--(b)--(c)--cycle;\node[left] at (a){$a$};\node[right] at (b){$b$};\node[above] at (c){$c$};\node[below] at (1.5,.1){$d(a,b)$};\node[left] at (.6,.9){$d(a,c)$};\node[right] at (2.2,1){$d(c,b)$};

\end{tikzpicture}
\figcap{解图：三个态向量在其仿射实平面中的距离三角形，距离由Hilbert范数定义。}\label{advA6:quality:A6-quantum-mechanics:a6}
\end{minipage}\end{center}

令 $x=a-c$、$y=c-b$，则 $a-b=x+y$。利用上题，
\[
\begin{aligned}
\|x+y\|^2
&=\|x\|^2+\|y\|^2+2\operatorname{Re}\langle x|y\rangle\\
&\le\|x\|^2+\|y\|^2+2\|x\|\|y\|
=(\|x\|+\|y\|)^2.
\end{aligned}
\]
两边取非负平方根即得
$\|a-b\|\le\|a-c\|+\|c-b\|$。

图~\ref{advA6:quality:A6-quantum-mechanics:a6} 中，沿 $a\to c\to b$ 的向量位移相加等于 $a\to b$，展开范数平方后用Cauchy--Schwarz控制交叉项，才得到长度的三角不等式。
\end{advsolution}

\begin{advdefinition}{共射线的归一化态矢}\label{advA6:def:same-ray}
若两个归一化态矢仅差一个全局相位，
\[
 |\psi'\rangle=\e^{\ii\theta}|\psi\rangle,
\]
则称它们在同一射线上。任何转移概率和可观测量期望都不依赖该全局相位。
\end{advdefinition}

\begin{advdefinition}{态等价与纯量子态}\label{advA6:def:state-equivalence}
任意非零共线态矢
$|\psi'\rangle=\lambda\e^{\ii\theta}|\psi\rangle$（$\lambda>0$）
代表同一纯态。纯量子态是非零 Hilbert 向量在此等价关系下的等价类；
计算时通常选取范数为 $1$ 的代表元。一般混态需用密度算符描述，不能都表示为单条态矢射线。
\end{advdefinition}

\begin{advdefinition}{左矢、右矢与对偶}\label{advA6:def:bra-ket}
空间中的向量记为右矢 $|a\rangle$。由 Riesz 表示定理，它唯一对应连续线性泛函
\[
 \langle a|:\ |\psi\rangle\mapsto\langle a|\psi\rangle,
\]
称为左矢。映射 $|a\rangle\mapsto\langle a|$ 是反线性的：
\[
 (\lambda|a\rangle+\mu|b\rangle)^\dagger
 =\lambda^*\langle a|+\mu^*\langle b|.
\]
\end{advdefinition}

\begin{advdefinition}{线性算符与反线性算符}\label{advA6:def:operator}
线性算符满足
\[
 A(\lambda|a\rangle+\mu|b\rangle)
 =\lambda A|a\rangle+\mu A|b\rangle.
\]
反线性算符则把系数取复共轭。时间反演是典型反线性对称操作；
一般可观测量由线性算符描述。
\end{advdefinition}

\begin{advexamplebox}{反线性算符诱导的对偶作用}\label{advA6:ex:antioperator}
设 $\Theta$ 是反线性算符。证明
\[
 \langle\Theta a|\Theta b\rangle
\]
对 $|a\rangle$ 线性、对 $|b\rangle$ 反线性，并说明反幺正条件如何保持概率。

\begin{center}\begin{minipage}{.96\linewidth}\centering
\begin{tikzpicture}[>=Latex,font=\small,scale=.85,every node/.style={fill=none}]
\draw[->] (-.4,0)--(3,0) node[right]{$x$};\draw[->] (0,-.4)--(0,2) node[above]{$y$};
\draw[DeepBlue,->] (0,0)--(1.3,1.1) node[above]{$\lambda$};\draw[dashed] (1.3,-1.1)--(1.3,1.1);\node[below] at (1.3,-1.1){$\lambda^*$};\node[right] at (2,1.5){复系数平面};

\end{tikzpicture}
\figcap{题图：反线性作用会共轭复系数，图只表示复系数而非Hilbert空间的物理位置。}\label{advA6:quality:A6-quantum-mechanics:q7}
\end{minipage}\end{center}
\end{advexamplebox}
\global\AdvHideHeadingtrue
\begin{advsolution}
\begin{center}\begin{minipage}{.96\linewidth}\centering
% EMAdvFigureHeading
\raggedright{\bfseries\color{DeepBlue}解答}\par\smallskip\centering

\begin{tikzpicture}[>=Latex,font=\small,scale=.85,every node/.style={fill=none}]
\draw[->] (-.4,0)--(3,0) node[right]{$x$};\draw[->] (0,-.4)--(0,2) node[above]{$y$};
\draw[DeepBlue,->] (0,0)--(1.3,1.1) node[above]{$\lambda$};\draw[dashed] (1.3,-1.1)--(1.3,1.1);\node[below] at (1.3,-1.1){$\lambda^*$};\node[right] at (2,1.5){复系数平面};

\end{tikzpicture}
\figcap{解图：反线性作用会共轭复系数，图只表示复系数而非Hilbert空间的物理位置。}\label{advA6:quality:A6-quantum-mechanics:a7}
\end{minipage}\end{center}

由反线性，
$\Theta(\lambda a+\mu b)=\lambda^*\Theta a+\mu^*\Theta b$。
因此
\[
 \langle\Theta(\lambda a+\mu b)|\Theta c\rangle
 =\lambda\langle\Theta a|\Theta c\rangle
 +\mu\langle\Theta b|\Theta c\rangle,
\]
而
\[
 \langle\Theta c|\Theta(\lambda a+\mu b)\rangle
 =\lambda^*\langle\Theta c|\Theta a\rangle
 +\mu^*\langle\Theta c|\Theta b\rangle.
\]
若
\[
 \langle\Theta a|\Theta b\rangle
 =\langle a|b\rangle^*=\langle b|a\rangle,
\]
若 $\Theta$ 还把整个空间映满，则称为反幺正；有限维等维情形下满射性由上述保内积条件推出。此时
\[
 |\langle\Theta a|\Theta b\rangle|^2
 =|\langle a|b\rangle|^2,
\]
所有转移概率保持不变。这里不能把反线性算符的“伴随”当作普通矩阵转置，
必须同时追踪复共轭的作用位置。

图~\ref{advA6:quality:A6-quantum-mechanics:a7} 中，$\lambda$ 与 $\lambda^*$ 在复系数平面关于实轴对称，故 $\Theta(\lambda|a\rangle)=\lambda^*\Theta|a\rangle$；再结合内积第一槽反线性，才能逐槽判断题中双线性类型。
\end{advsolution}

\begin{advdefinition}{伴随算符}\label{advA6:def:adjoint}
稠密定义的线性算符 $A$ 的伴随 $A^\dagger$ 由
\[
 \langle\psi|A\phi\rangle
 =\langle A^\dagger\psi|\phi\rangle
\]
定义，其中还必须指定两者的定义域。有限维中，
$[A^\dagger]=[A]^{*\,\mathsf T}$。
\end{advdefinition}

\begin{advdefinition}{可观测量}\label{advA6:def:observable}
可观测量由自伴算符表示。自伴不仅要求形式上
$A^\dagger=A$，还要求定义域相同。
其谱为实数，并由谱投影给出完备分解；有限维时即 Hermitian 矩阵。
\end{advdefinition}

\subsection{位置表象、动量表象和常用算符}

位置与动量广义本征态满足
\[
 \langle x'|x\rangle=\delta(x-x'),\qquad
 \langle p'|p\rangle=\delta(p-p'),
\]
\[
 \langle x|p\rangle=
 \frac1{\sqrt{2\pi\hbar}}\e^{\ii px/\hbar}.
\]
完备性为
\[
 \int|x\rangle\langle x|\,\dd x=I,\qquad
 \int|p\rangle\langle p|\,\dd p=I.
\]
在位置表象中
\[
 \widehat x=x,\qquad
 \widehat p=-\ii\hbar\frac{\dd}{\dd x},\qquad
 [\widehat x,\widehat p]=\ii\hbar.
\]
平移算符
\[
 T(\varepsilon)=\exp\!\left(-\frac{\ii}{\hbar}
 \varepsilon\widehat p\right)
\]
作用为
$(T(\varepsilon)\psi)(x)=\psi(x-\varepsilon)$；
正负号取决于是主动平移态还是被动平移坐标，必须保持约定一致。

\begin{advexamplebox}{投影、位置、动量和平移算符的伴随性质}\label{advA6:ex:operator-adjoints}
证明位置/动量谱投影、位置算符和动量算符为自伴对象，
而非零平移算符一般不是 Hermitian，却是幺正算符，且
$T(\varepsilon)^\dagger=T(-\varepsilon)$。

\begin{center}\begin{minipage}{.96\linewidth}\centering
\begin{tikzpicture}[>=Latex,font=\small,scale=.85,every node/.style={fill=none}]

\draw[->] (-1,0)--(4.5,0) node[right]{$x$};\draw[DeepBlue,domain=-.8:2,samples=60] plot(\x,{exp(-3*(\x-.5)^2)});\draw[DeepBlue,dashed,domain=.5:3.5,samples=60] plot(\x,{exp(-3*(\x-2)^2)});\draw[->] (.5,1.3)--(2,1.3) node[midway,above]{$+\epsilon$};\node[below] at (2,-.5){虚线为 $\psi(x-\epsilon)$};

\end{tikzpicture}
\figcap{题图：实轴上的波包与主动平移后的波包，图中取 $\epsilon>0$。}\label{advA6:quality:A6-quantum-mechanics:q8}
\end{minipage}\end{center}
\end{advexamplebox}
\global\AdvHideHeadingtrue
\begin{advsolution}
\begin{center}\begin{minipage}{.96\linewidth}\centering
% EMAdvFigureHeading
\raggedright{\bfseries\color{DeepBlue}解答}\par\smallskip\centering

\begin{tikzpicture}[>=Latex,font=\small,scale=.85,every node/.style={fill=none}]

\draw[->] (-1,0)--(4.5,0) node[right]{$x$};\draw[DeepBlue,domain=-.8:2,samples=60] plot(\x,{exp(-3*(\x-.5)^2)});\draw[DeepBlue,dashed,domain=.5:3.5,samples=60] plot(\x,{exp(-3*(\x-2)^2)});\draw[->] (.5,1.3)--(2,1.3) node[midway,above]{$+\epsilon$};\node[below] at (2,-.5){虚线为 $\psi(x-\epsilon)$};

\end{tikzpicture}
\figcap{解图：实轴上的波包与主动平移后的波包，图中取 $\epsilon>0$。}\label{advA6:quality:A6-quantum-mechanics:a8}
\end{minipage}\end{center}

先用有限测量区间 $\Delta$ 写位置谱投影：
\[
 (P_x(\Delta)\psi)(x)=\mathbf1_\Delta(x)\psi(x),\qquad
 P_x(\Delta)^\dagger=P_x(\Delta),\quad P_x(\Delta)^2=P_x(\Delta).
\]
动量谱投影在动量表象中有同样的乘法形式，或由 Fourier 酉变换从该形式得到。
$|x\rangle\langle x|$ 是形式上的谱密度，不是单个正规的秩一投影。
$\widehat x$ 在 $D(\widehat x)=\{\psi\in L^2:x\psi\in L^2\}$ 上是自伴乘法算符。对动量算符，
\[
 \langle\phi|\widehat p\psi\rangle
 =-\ii\hbar\int\phi^*\psi'\,\dd x
 =-\ii\hbar[\phi^*\psi]_{\rm bdry}
 +\ii\hbar\int(\phi^*)'\psi\,\dd x.
\]
在波函数于无穷远消失或采用相容边界条件时，边界项为零，得到
$\langle\phi|\widehat p\psi\rangle
=\langle\widehat p\phi|\psi\rangle$；
在整条实线上取 $D(\widehat p)=H^1(\mathbb R)$，可核验伴随定义域相同，因而自伴。有限区间时仅消掉边界项只能证明对称，还须另行指定自伴边界条件；例如两端都取Dirichlet条件的动量算符并不自动自伴。

因为 $\widehat p$ 自伴，
\[
 T(\varepsilon)^\dagger
 =\exp\!\left(+\frac{\ii\varepsilon\widehat p}{\hbar}\right)
 =T(-\varepsilon)=T(\varepsilon)^{-1}.
\]
所以 $T$ 幺正。除 $\varepsilon=0$ 或特殊离散谱情形外，
$T^\dagger\ne T$，本征值 $\e^{-\ii\varepsilon p/\hbar}$ 也不是实数，
故它本身不是可观测量。

图~\ref{advA6:quality:A6-quantum-mechanics:a8} 中，主动向右平移作用在函数上是自变量替换 $x\mapsto x-\epsilon$，其逆为 $-\epsilon$；内积保持说明幺正，不能由此直接断言平移算符自伴。
\end{advsolution}

\begin{advexamplebox}{位置、动量与平移的广义本征态}\label{advA6:ex:generalized-eigenstates}
说明位置算符、动量算符和平移算符的本征态与完备性；
解释连续谱密度、区间投影与位置算符三者的区别。

\begin{center}\begin{minipage}{.96\linewidth}\centering
\begin{tikzpicture}[>=Latex,font=\small,scale=.85,every node/.style={fill=none}]

\draw[->] (-1,0)--(4.7,0) node[right]{$x$};\draw[very thick,DeepBlue] (.5,0)--(3,0);\draw (.5,-.15)--(.5,.15) (3,-.15)--(3,.15);\node[below] at (.5,0){$a$};\node[below] at (3,0){$b$};\draw[<->] (1.5,.7)--(1.9,.7) node[midway,above]{$\mathrm dx$};\node[above] at (3.2,1.3){位置谱区间投影};

\end{tikzpicture}
\figcap{题图：连续位置谱中的一个区间及微元，区间投影与点本征态不同。}\label{advA6:quality:A6-quantum-mechanics:q9}
\end{minipage}\end{center}
\end{advexamplebox}
\global\AdvHideHeadingtrue
\begin{advsolution}
\begin{center}\begin{minipage}{.96\linewidth}\centering
% EMAdvFigureHeading
\raggedright{\bfseries\color{DeepBlue}解答}\par\smallskip\centering

\begin{tikzpicture}[>=Latex,font=\small,scale=.85,every node/.style={fill=none}]

\draw[->] (-1,0)--(4.7,0) node[right]{$x$};\draw[very thick,DeepBlue] (.5,0)--(3,0);\draw (.5,-.15)--(.5,.15) (3,-.15)--(3,.15);\node[below] at (.5,0){$a$};\node[below] at (3,0){$b$};\draw[<->] (1.5,.7)--(1.9,.7) node[midway,above]{$\mathrm dx$};\node[above] at (3.2,1.3){位置谱区间投影};

\end{tikzpicture}
\figcap{解图：连续位置谱中的一个区间及微元，区间投影与点本征态不同。}\label{advA6:quality:A6-quantum-mechanics:a9}
\end{minipage}\end{center}

\[
 \widehat x|x_0\rangle=x_0|x_0\rangle,\qquad
 \widehat p|p_0\rangle=p_0|p_0\rangle.
\]
$|x_0\rangle$ 和 $|p_0\rangle$ 不是 Hilbert 空间中的可归一向量，
而是 rigged Hilbert space 中的广义向量；其 $\delta$ 归一化与积分完备性如上。
任意态均可展开
\[
 |\psi\rangle=\int\psi(x)|x\rangle\,\dd x
 =\int\tilde\psi(p)|p\rangle\,\dd p.
\]
平移算符与动量算符具有同一组广义本征态：
\[
 T(\varepsilon)|p\rangle
 =\e^{-\ii\varepsilon p/\hbar}|p\rangle.
\]
严格位置测量由投影值测度
\[
 P_x(\Delta)=\int_\Delta|x\rangle\langle x|\,\dd x
\]
描述。非平凡区间的投影有 $0,1$ 两个谱值，回答“是否在这个区间”；
位置算符则是 $\widehat x=\int x\,\dd P_x(x)$，给出连续的位置值。
在 $L^2(\mathbb R)$ 的通常位置表象中，单点集合测度为零，所以
$P_x(\{x_0\})=0$。形式符号 $|x_0\rangle\langle x_0|$ 不能直接平方成自身，
否则会遇到无意义的 $\delta(0)$；它不是具有 $0,1$ 两个谱值的普通秩一投影。

图~\ref{advA6:quality:A6-quantum-mechanics:a9} 中，蓝色区间对应 $P_{[a,b]}=\int_a^b|x\rangle\langle x|\,\mathrm dx$；点核是算符值密度，位置算符还需再乘坐标 $x$ 积分，三者的量纲和作用都不同。
\end{advsolution}

\subsection{有限维态、矩阵与自旋}

在二维自旋空间中取基
$|1\rangle=|\uparrow\rangle$、
$|2\rangle=|\downarrow\rangle$，
态写为
\[
 |\psi\rangle=
 \psi_1|\uparrow\rangle+\psi_2|\downarrow\rangle,
\qquad
 |\psi_1|^2+|\psi_2|^2=1.
\]

\begin{advexamplebox}{算符矩阵与伴随矩阵}\label{advA6:ex:matrix-adjoint}
在正交归一基中证明
\[
 [A^\dagger]=[A]^{*\,\mathsf T}.
\]
并说明 Hermitian 矩阵的元素条件。

\begin{center}\begin{minipage}{.96\linewidth}\centering
\begin{tikzpicture}[>=Latex,font=\small,scale=.85,every node/.style={fill=none}]

\fill[DeepBlue!12] (.8,1.2) rectangle (1.6,1.8);\draw (0,0) grid[xstep=.8,ystep=.6] (2.4,1.8);\node at (1.2,1.5){$A_{ij}$};\node[above] at (1.2,2){正交归一基};\node[below] at (1.2,-.4){行 $i$，列 $j$（以非对角元示意）};

\end{tikzpicture}
\figcap{题图：算符矩阵的一个行列元素，索引由所选正交归一基决定。}\label{advA6:quality:A6-quantum-mechanics:q10}
\end{minipage}\end{center}
\end{advexamplebox}
\global\AdvHideHeadingtrue
\begin{advsolution}
\begin{center}\begin{minipage}{.96\linewidth}\centering
% EMAdvFigureHeading
\raggedright{\bfseries\color{DeepBlue}解答}\par\smallskip\centering

\begin{tikzpicture}[>=Latex,font=\small,scale=.85,every node/.style={fill=none}]

\fill[DeepBlue!12] (.8,1.2) rectangle (1.6,1.8);\draw (0,0) grid[xstep=.8,ystep=.6] (2.4,1.8);\node at (1.2,1.5){$A_{ij}$};\node[above] at (1.2,2){$[A]$};
\begin{scope}[xshift=4.2cm]\fill[DeepBlue!12] (0,.6) rectangle (.8,1.2);\draw (0,0) grid[xstep=.8,ystep=.6] (2.4,1.8);\node at (.4,.9){$A_{ij}^*$};\node[above] at (1.2,2){$[A^\dagger]$};\end{scope}
\node[below] at (3.3,-.5){原来的 $(i,j)$ 元素进入伴随矩阵的 $(j,i)$ 位置并取复共轭};

\end{tikzpicture}
\figcap{解图：非对角矩阵元的位置对换与复共轭同时进行。}\label{advA6:quality:A6-quantum-mechanics:a10}
\end{minipage}\end{center}

矩阵元 $A_{ij}=\langle i|A|j\rangle$。由伴随定义，
\[
 (A^\dagger)_{ij}
 =\langle i|A^\dagger|j\rangle
 =\langle j|A|i\rangle^*
 =A_{ji}^*.
\]
所以 $[A^\dagger]=[A]^{*\,\mathsf T}$。
若 $A=A^\dagger$，则
$A_{ij}=A_{ji}^*$，对角元 $A_{ii}$ 必为实数；
反之满足这些元素关系的矩阵就是 Hermitian 矩阵。
投影和 Pauli 矩阵均满足该条件；升、降算符互为伴随。

图~\ref{advA6:quality:A6-quantum-mechanics:a10} 中，行列交换来自交换内积的两个基向量，复共轭来自伴随定义；因此必须把转置和共轭同时执行，不能仅对每个元素取共轭而保持索引不变。
\end{advsolution}

Hermitian 算符的本征值实、不同本征值的本征态正交。
若 $A|a\rangle=a|a\rangle$，
\[
 a\langle a|a\rangle
 =\langle a|A|a\rangle
 =\langle a|A^\dagger|a\rangle
 =a^*\langle a|a\rangle,
\]
故 $a=a^*$。对 $a\ne b$，
$(a-b)\langle b|a\rangle=0$，故正交。
离散谱的完备性写成
$\sum_{i,j}|\psi_{ij}\rangle\langle\psi_{ij}|=I$；
连续谱则以积分和 Dirac $\delta$ 取代求和与 Kronecker $\delta$。

\begin{advtheorem}{投影测量规则}\label{advA6:thm:copenhagen}
设可观测量具有谱分解
\[
 A=\sum_iA_iP_i,
\]
其中 $P_i$ 是本征值 $A_i$ 的简并本征空间投影。归一化态
$|\psi\rangle$ 测得 $A_i$ 的概率为
\[
 \boxed{p_i=\langle\psi|P_i|\psi\rangle}.
\]
在理想选择性测量中，对 $p_i>0$ 的实际结果，得到 $A_i$ 后的态为
\[
 \boxed{|\psi_i'\rangle=
 \frac{P_i|\psi\rangle}{\sqrt{p_i}}}.
\]
连续谱用投影值测度替代离散求和。
\end{advtheorem}

\begin{advdefinition}{期望值、矩阵元与跃迁振幅}\label{advA6:def:expectation-transition}
归一化态中算符 $A$ 的期望值为 $\langle A\rangle_\psi=\langle\psi|A|\psi\rangle$。
矩阵元 $A_{\phi\psi}=\langle\phi|A|\psi\rangle$ 与 $A$ 具有相同量纲，
不能无条件把它的模方当作跃迁概率。真正的演化振幅是
\[
 \mathcal A_{\psi\to\phi}(t)=\langle\phi|U(t,0)|\psi\rangle,
\]
其中 $U$ 是时间演化算符。对时间无关的 $H_0$ 的不同本征态 $i,f$，
若附加弱扰动 $V(t)$，相互作用表象的一阶振幅为
\[
 c_f^{(1)}(t)=-\frac{\ii}{\hbar}\int_0^t
 \langle f|V(t')|i\rangle\,
 \e^{\ii(E_f-E_i)t'/\hbar}\,\dd t'.
\]
积分和 $1/\hbar$ 使振幅无量纲；在一阶微扰适用且该一阶振幅非零时，$|c_f^{(1)}|^2$ 给出相应跃迁概率的最低非零阶。若一阶振幅被选择定则禁止，则需继续计算高阶项。
\end{advdefinition}

\begin{advdefinition}{对易子}\label{advA6:def:commutator}
两个算符的对易子为
\[
 [A,B]=AB-BA.
\]
$[A,B]=0$ 时称两者对易或相容。对易关系需在共同稠密定义域上理解；
仅有形式微分表达式相等还不够。
\end{advdefinition}

\begin{advtheorem}{对易自伴算符的共同本征分解}\label{advA6:thm:commuting-observables}
有限维中，两个对易 Hermitian 算符可被同一幺正矩阵同时对角化，
因而存在共同正交归一本征基。无限维中，强对易的自伴算符具有共同谱测度。
\end{advtheorem}

\noindent\textit{证明（有限维）。}
取 $A$ 的某一本征空间 $V_a$。若 $|x\rangle\in V_a$，
\[
 A(B|x\rangle)=B(A|x\rangle)=aB|x\rangle,
\]
故 $B$ 把 $V_a$ 映入自身。在每个 $V_a$ 中限制的 $B$ 仍为 Hermitian，
可选取其正交归一本征基。对所有 $A$ 的本征空间合并这些基，
便得到 $A,B$ 的共同正交归一本征基。
一组对易可观测量若共同量子数能唯一标记每个共同本征态，称为对易可观测量完全集。

\begin{advdefinition}{Pauli 矩阵}\label{advA6:def:pauli}
\[
 \sigma_x=\begin{pmatrix}0&1\\1&0\end{pmatrix},\quad
 \sigma_y=\begin{pmatrix}0&-\ii\\\ii&0\end{pmatrix},\quad
 \sigma_z=\begin{pmatrix}1&0\\0&-1\end{pmatrix}.
\]
它们满足
\[
 \sigma_i\sigma_j=\delta_{ij}I+\ii\varepsilon_{ijk}\sigma_k,
\qquad
 [\sigma_i,\sigma_j]=2\ii\varepsilon_{ijk}\sigma_k.
\]
自旋算符为 $\vect S=(\hbar/2)\vect\sigma$。
\end{advdefinition}

\begin{advexamplebox}{Pauli 本征态、投影与升降算符}\label{advA6:ex:pauli-eigenstates}
\begin{enumerate}
\item 求 $P_\uparrow=|\uparrow\rangle\langle\uparrow|$、
$P_\downarrow$ 的本征值和本征空间；
\item 求 $\sigma_x,\sigma_y,\sigma_z$ 的本征态；
\item 定义
$\sigma_\pm=(\sigma_x\pm\ii\sigma_y)/2$，
说明其作用并证明没有非零本征值。
\end{enumerate}

\begin{center}\begin{minipage}{.96\linewidth}\centering
\begin{tikzpicture}[>=Latex,font=\small,scale=.85,every node/.style={fill=none}]

\draw[->] (0,-1.4)--(0,1.7) node[above]{$z$};\draw[->] (-1.5,0)--(1.7,0) node[right]{$x$};\draw[->] (-.9,-.7)--(.9,.7) node[above right]{$y$};\draw[dashed] (0,0) circle (1.1);\fill[DeepBlue] (0,1.1) circle (2pt) node[left]{$|\uparrow\rangle$};\fill[DeepBlue] (0,-1.1) circle (2pt) node[left]{$|\downarrow\rangle$};

\end{tikzpicture}
\figcap{题图：自旋1/2态的Bloch球示意，南北极对应 $\sigma_z$ 的两个本征方向。}\label{advA6:quality:A6-quantum-mechanics:q11}
\end{minipage}\end{center}
\end{advexamplebox}
\global\AdvHideHeadingtrue
\begin{advsolution}
\begin{center}\begin{minipage}{.96\linewidth}\centering
% EMAdvFigureHeading
\raggedright{\bfseries\color{DeepBlue}解答}\par\smallskip\centering

\begin{tikzpicture}[>=Latex,font=\small,scale=.85,every node/.style={fill=none}]

\draw[->] (0,-1.4)--(0,1.7) node[above]{$z$};\draw[->] (-1.5,0)--(1.7,0) node[right]{$x$};\draw[->] (-.9,-.7)--(.9,.7) node[above right]{$y$};\draw[dashed] (0,0) circle (1.1);\fill[DeepBlue] (0,1.1) circle (2pt) node[left]{$|\uparrow\rangle$};\fill[DeepBlue] (0,-1.1) circle (2pt) node[left]{$|\downarrow\rangle$};

\end{tikzpicture}
\figcap{解图：自旋1/2态的Bloch球示意，南北极对应 $\sigma_z$ 的两个本征方向。}\label{advA6:quality:A6-quantum-mechanics:a11}
\end{minipage}\end{center}

\[
 P_\uparrow=\begin{pmatrix}1&0\\0&0\end{pmatrix},\qquad
 P_\downarrow=\begin{pmatrix}0&0\\0&1\end{pmatrix}.
\]
两者本征值均为 $0,1$，各自的完整本征基是
$|\uparrow\rangle,|\downarrow\rangle$；
投影本身可作为二值可观测量，但不是“位置或自旋分量”本身。

三个 Pauli 矩阵的 $\pm1$ 本征态可取
\[
 |x\pm\rangle=\frac1{\sqrt2}
 (|\uparrow\rangle\pm|\downarrow\rangle),
\]
\[
 |y\pm\rangle=\frac1{\sqrt2}
 (|\uparrow\rangle\pm\ii|\downarrow\rangle),
\qquad
 |z+\rangle=|\uparrow\rangle,\quad
 |z-\rangle=|\downarrow\rangle.
\]
每组正交完备。

\[
 \sigma_+=
 \begin{pmatrix}0&1\\0&0\end{pmatrix},\qquad
 \sigma_-=
 \begin{pmatrix}0&0\\1&0\end{pmatrix}
\]
（若把基矢次序反置，矩阵位置相应互换）。它们把一个 $z$ 本征态映到另一个，
并把最高/最低态映为零。由于
$\sigma_\pm^2=0$，任一本征值 $\lambda$ 必满足 $\lambda^2=0$，
故只有 $\lambda=0$；核空间虽有非零向量，但不存在非零本征值。

图~\ref{advA6:quality:A6-quantum-mechanics:a11} 中，投影算符保留各自极点的态并湮灭正交态；$\sigma_+$ 从下态升到上态而在上态作用为零，$\sigma_-$ 相反。Bloch方向不等于一个粒子的空间轨道。
\end{advsolution}

\begin{advexamplebox}{Pauli 矩阵指数与自旋旋转}\label{advA6:ex:pauli-exponential}
对单位矢量 $\vect n$，求
\[
 U=\exp(\ii\theta\,\vect n\cdot\vect\sigma).
\]

\begin{center}\begin{minipage}{.96\linewidth}\centering
\begin{tikzpicture}[>=Latex,font=\small,scale=.85,every node/.style={fill=none}]

\draw[->] (0,-1.4)--(0,1.7) node[above]{$z$};\draw[->] (-1.5,0)--(1.7,0) node[right]{$x$};\draw[->] (-.9,-.7)--(.9,.7) node[above right]{$y$};\draw[dashed] (0,0) circle (1.1);\fill[DeepBlue] (0,1.1) circle (2pt) node[left]{$|\uparrow\rangle$};\fill[DeepBlue] (0,-1.1) circle (2pt) node[left]{$|\downarrow\rangle$};
\draw[DeepBlue,->,thick] (0,0)--(1,1.3) node[right]{$\boldsymbol n$};
\end{tikzpicture}
\figcap{题图：单位方向 $\boldsymbol n$ 决定自旋算符指数的轴，指数参数为 $\theta$。}\label{advA6:quality:A6-quantum-mechanics:q12}
\end{minipage}\end{center}
\end{advexamplebox}
\global\AdvHideHeadingtrue
\begin{advsolution}
\begin{center}\begin{minipage}{.96\linewidth}\centering
% EMAdvFigureHeading
\raggedright{\bfseries\color{DeepBlue}解答}\par\smallskip\centering

\begin{tikzpicture}[>=Latex,font=\small,scale=.85,every node/.style={fill=none}]

\draw[->] (0,-1.4)--(0,1.7) node[above]{$z$};\draw[->] (-1.5,0)--(1.7,0) node[right]{$x$};\draw[->] (-.9,-.7)--(.9,.7) node[above right]{$y$};\draw[dashed] (0,0) circle (1.1);\fill[DeepBlue] (0,1.1) circle (2pt) node[left]{$|\uparrow\rangle$};\fill[DeepBlue] (0,-1.1) circle (2pt) node[left]{$|\downarrow\rangle$};
\draw[DeepBlue,->,thick] (0,0)--(1,1.3) node[right]{$\boldsymbol n$};
\end{tikzpicture}
\figcap{解图：单位方向 $\boldsymbol n$ 决定自旋算符指数的轴，指数参数为 $\theta$。}\label{advA6:quality:A6-quantum-mechanics:a12}
\end{minipage}\end{center}

利用 Pauli 乘法，
\[
 (\vect n\cdot\vect\sigma)^2
 =n_in_j(\delta_{ij}I+\ii\varepsilon_{ijk}\sigma_k)
 =|\vect n|^2I=I.
\]
指数级数的偶次幂与奇次幂分别求和：
\[
\begin{aligned}
 U&=\sum_{m=0}^\infty
 \frac{(\ii\theta)^{2m}}{(2m)!}I
 +\sum_{m=0}^\infty
 \frac{(\ii\theta)^{2m+1}}{(2m+1)!}
 \vect n\cdot\vect\sigma\\
 &=\boxed{\cos\theta\,I+
 \ii\sin\theta\,\vect n\cdot\vect\sigma}.
\end{aligned}
\]
它满足 $U^\dagger U=I$、$\det U=1$。
物理空间把自旋旋转角 $\varphi$ 时，旋转算符为
\[
 U_{\rm rot}(\varphi)=
 \exp\!\left(-\frac{\ii\varphi}{2}
 \vect n\cdot\vect\sigma\right).
\]
因半角出现，自旋子旋转 $2\pi$ 后变号，旋转 $4\pi$ 才回到原态矢；
但全局负号不改变单次态的物理射线。

图~\ref{advA6:quality:A6-quantum-mechanics:a12} 中，沿 $\boldsymbol n$ 的两个本征态分别获得 $e^{\pm\mathrm i\theta}$ 相位；相对相位变为 $2\theta$，对 $U=e^{+\mathrm i\theta\boldsymbol n\cdot\boldsymbol\sigma}$，Bloch向量的主动转角为 $-2\theta$。
\end{advsolution}

\subsubsection*{Bloch 矢量上的旋转及可观测检验}
任意纯自旋态的密度矩阵都可写成
\[
 \rho=|\psi\rangle\langle\psi|
 =\frac12\bigl(I+\vect r\cdot\vect\sigma\bigr),
 \qquad |\vect r|=1.
\]
混态仍具有同一形式，只是 $|\vect r|<1$。对空间转角 $\varphi$，令
$U=\exp[-\ii\varphi\vect n\cdot\vect\sigma/2]$，则
\[
 \rho'=U\rho U^\dagger
 =\frac12\bigl(I+\vect r'\cdot\vect\sigma\bigr),
\]
其中直接展开 Pauli 乘积可得 Rodrigues 公式
\[
 \boxed{\vect r'
 =\vect r\cos\varphi
 +(\vect n\times\vect r)\sin\varphi
 +\vect n(\vect n\cdot\vect r)(1-\cos\varphi)}.
\]
这说明自旋态矢使用半角旋转，而所有可观测的 Bloch 矢量仍按真实空间角
$\varphi$ 旋转。沿单位方向 $\vect a$ 测量自旋，得到正结果的概率为
\[
 P_{\vect a,+}
 =\operatorname{tr}\!\left[
 \rho\,\frac{I+\vect a\cdot\vect\sigma}{2}\right]
 =\frac{1+\vect a\cdot\vect r}{2}.
\]
因此 $2\pi$ 旋转产生的整体负号在单一路径概率中消失；若把旋转过的态与
未旋转的参考路径相干叠加，相对负号则可转化为干涉条纹的半周期移动。
